\section{Objetivos}

\subsection{Objetivo general}
Desarrollar un prototipo web que apoye la comprensión de los contenidos de la asignatura Fundamentos de Interpretación y Compilación de Lenguajes de Programación (LP) de la Universidad del Valle mediante la visualización de estructuras internas y el uso de lenguajes orientados por sintaxis.

\subsection{Objetivos específicos y resultados esperados}
\begin{table}[H]
    \centering
    \begin{tabular}{|p{11cm}|p{5cm}|}
        \hline
        \textbf{Objetivo específico} & \textbf{Sección del documento} \\
        \hline
        Identificar las dificultades conceptuales que enfrentan los estudiantes de la asignatura, con el propósito de orientar la selección de ejemplos, casos de uso y mecanismos de visualización.&Capítulo 3\\
        \hline
        Caracterizar los fundamentos teóricos y técnicos relacionados con las gramáticas formales, el análisis sintáctico y la generación de árboles de sintaxis abstracta, que servirán como base conceptual para el diseño del prototipo.&Capítulo 2 subsecciones 2.2.1 y 2.2.4 y Capítulo 4 sección 4.1 \\
        \hline
        Diseñar la arquitectura web y las interfaces de usuario, integrando estrategias de visualización que favorezcan la comprensión de las estructuras sintácticas.& Capítulo 4\\
        \hline
        Implementar un prototipo funcional que permita definir gramáticas, visualizar los árboles de sintaxis abstracta y los cambios en el estado durante la evaluación de los programas escritos por el usuario.&Capítulo 5\\
        \hline
        Evaluar la usabilidad y funcionalidad del prototipo mediante pruebas con estudiantes inscritos en la asignatura y con el docente responsable.&Capítulo 6\\
        \hline
    \end{tabular}
    \caption{Objetivos Específicos}
    \label{tab:objEspecificos}
\end{table}

